% P46 | 5.8 Profit and loss | Result: the buyer's profit and loss under the pricing measure
% guide: dev_notes/paragraph_guides/P46_result_the_buyer_s_profit_and_loss_under_the_pri.txt
The buyer pays the fair value $V(0,S_0)=\AmVzero$ and exercises at the first date on which the path enters the exercise region.
Paths are simulated in the tree with the pricing weights $(p,1-p)$, because under the pricing measure the real-world drift $\mu$ plays no role.
The mean profit is \PnLmean\ with a standard error of \PnLse, consistent with zero, as it must be by the definition of the fair value.
A vectorised run with $10^5$ paths gives \PnLQmean\ with a standard error of \PnLQse, which confirms it.
The standard deviation of the profit is \PnLsd, the mean stopping time is \PnLtau\ years and the probability of exercising before maturity is \PnLpex.
Roughly half of the paths are therefore stopped early.
The distribution in Figure~\ref{fig:3a} has a point mass at $-V_0$ for the paths that end out of the money and a long right tail, and Table~\ref{tbl:pnlgrid} repeats the experiment on the nine $(\sigma,r)$ combinations.
The spread grows with $\sigma$ and shrinks with $r$, the probability of early exercise rises with $\sigma$, and for $r=0$ the option is never exercised early, so that $\tau=T$.
